% GENERATED FILE. Edit canonical lessons, then run npm run generate:books.
% canonical-source-hash: fnv1a64:a804d955c56395f0
% canonical-lessons: FR-C504-musculation, FR-C504-haltere, FR-C504-etirer, FR-C504-sueur, FR-C504-echauffer

\chapter{À la salle}
\label{ch:fr-a-la-salle}
\hlchaptermodality{504}

\begin{chapteropening}
\textbf{By the end of this chapter:} \emph{I can talk about working out at the gym.}

Complete the last lesson of chapter 504: I can talk about working out at the gym.
\end{chapteropening}

\section[la musculation]{la musculation — weight training, bodybuilding}

\label{lesson:FR-C504-musculation}



\begin{quote}
Before the new one: say the French for to water, then the French for watering can.
\end{quote}

\subsection*{You'll want to know: la musculation}
\textbf{la musculation} — \textquotedblleft{}weight training, bodybuilding\textquotedblright{}.

\begin{quote}
Je fais de la musculation deux fois par semaine.
\end{quote}

\begin{quote}
\emph{I do weight training twice a week.}
\end{quote}

Informally \emph{la muscu}.

\begin{cousinweb}
\emph{la musculation} is built on \emph{muscle}, a learned borrowing from Latin \emph{musculus} \textquotedblleft{}a little mouse, a muscle\textquotedblright{}, from \emph{mūs} \textquotedblleft{}mouse\textquotedblright{} (Lewis \& Short): a moving muscle looked like a mouse under the skin.
\end{cousinweb}

\begin{grammarlens}[title={What you've built}]
One more word for telling, explaining and describing at B1.
\end{grammarlens}

\subsection*{Your turn}
\begin{itemize}
\raggedright
  \item \emph{Say it:} \emph{la musculation}
  \item \emph{Say it:} \emph{la musculation}, once more
  \item \emph{Say it:} a sentence of your own with \emph{la musculation}
  \item \emph{Recall:} say the French for to water, then the French for watering can, then say \emph{la musculation} again
\end{itemize}

\begin{tcolorbox}[breakable,colback=teal!4,colframe=teal!35!black,title={Before you move on}]
What does la musculation mean? (\textquotedblleft{}Weight training, bodybuilding\textquotedblright{}.) Say it once more.
\end{tcolorbox}
\section[l'haltère]{l'haltère — dumbbell, weight}

\label{lesson:FR-C504-haltere}



\begin{quote}
Before the new one: say the French for watering can, then the French for weight training.
\end{quote}

\subsection*{You'll want to know: l'haltère}
\textbf{l'haltère} — \textquotedblleft{}dumbbell, weight\textquotedblright{}.

\begin{quote}
Il soulève des haltères devant le miroir.
\end{quote}

\begin{quote}
\emph{He lifts dumbbells in front of the mirror.}
\end{quote}

Masculine: \emph{un haltère}. The \emph{h} is silent.

\begin{cousinweb}
\emph{l'haltère} comes, through Latin, from Greek \emph{haltēres}, weights that athletes held while jumping, built on \emph{hállesthai} \textquotedblleft{}to jump\textquotedblright{}.
\end{cousinweb}

\begin{grammarlens}[title={What you've built}]
One more word for telling, explaining and describing at B1.
\end{grammarlens}

\subsection*{Your turn}
\begin{itemize}
\raggedright
  \item \emph{Say it:} \emph{l'haltère}
  \item \emph{Say it:} \emph{l'haltère}, once more
  \item \emph{Say it:} a sentence of your own with \emph{l'haltère}
  \item \emph{Recall:} say the French for watering can, then the French for weight training, then say \emph{l'haltère} again
\end{itemize}

\begin{tcolorbox}[breakable,colback=teal!4,colframe=teal!35!black,title={Before you move on}]
What does l'haltère mean? (\textquotedblleft{}Dumbbell, weight\textquotedblright{}.) Say it once more.
\end{tcolorbox}
\section[étirer]{étirer — to stretch; (s'étirer) to stretch (oneself)}

\label{lesson:FR-C504-etirer}



\begin{quote}
Before the new one: say the French for weight training, then the French for dumbbell.
\end{quote}

\subsection*{You'll want to know: étirer}
\textbf{étirer} — \textquotedblleft{}to stretch; (s'étirer) to stretch (oneself)\textquotedblright{}.

\begin{quote}
Après le sport, je m'étire pendant dix minutes.
\end{quote}

\begin{quote}
\emph{After sport I stretch for ten minutes.}
\end{quote}

Mostly reflexive for the body: \emph{s'étirer}. \emph{Un étirement} — a stretch.

\begin{cousinweb}
\emph{étirer} is built from \emph{é-} and \emph{tirer} \textquotedblleft{}to pull\textquotedblright{}. The origin of \emph{tirer} is not settled.
\end{cousinweb}

\begin{grammarlens}[title={What you've built}]
One more word for telling, explaining and describing at B1.
\end{grammarlens}

\subsection*{Your turn}
\begin{itemize}
\raggedright
  \item \emph{Say it:} \emph{étirer}
  \item \emph{Say it:} \emph{étirer}, once more
  \item \emph{Say it:} a sentence of your own with \emph{étirer}
  \item \emph{Recall:} say the French for weight training, then the French for dumbbell, then say \emph{étirer} again
\end{itemize}

\begin{tcolorbox}[breakable,colback=teal!4,colframe=teal!35!black,title={Before you move on}]
What does étirer mean? (\textquotedblleft{}To stretch; (s'étirer) to stretch (oneself)\textquotedblright{}.) Say it once more.
\end{tcolorbox}
\section[la sueur]{la sueur — sweat}

\label{lesson:FR-C504-sueur}



\begin{quote}
Before the new one: say the French for dumbbell, then the French for to stretch.
\end{quote}

\subsection*{You'll want to know: la sueur}
\textbf{la sueur} — \textquotedblleft{}sweat\textquotedblright{}.

\begin{quote}
Après une heure de course, il était tout en sueur.
\end{quote}

\begin{quote}
\emph{After an hour's running he was dripping with sweat.}
\end{quote}

\emph{En sueur} — sweating.

\begin{cousinweb}
\emph{la sueur} is inherited, through Old French \emph{suor}, from Latin \emph{sūdor} \textquotedblleft{}sweat\textquotedblright{} (EtymDB), related to \emph{sūdāre} \textquotedblleft{}to sweat\textquotedblright{} (Lewis \& Short).
\end{cousinweb}

\begin{grammarlens}[title={What you've built}]
One more word for telling, explaining and describing at B1.
\end{grammarlens}

\subsection*{Your turn}
\begin{itemize}
\raggedright
  \item \emph{Say it:} \emph{la sueur}
  \item \emph{Say it:} \emph{la sueur}, once more
  \item \emph{Say it:} a sentence of your own with \emph{la sueur}
  \item \emph{Recall:} say the French for dumbbell, then the French for to stretch, then say \emph{la sueur} again
\end{itemize}

\begin{tcolorbox}[breakable,colback=teal!4,colframe=teal!35!black,title={Before you move on}]
What does la sueur mean? (\textquotedblleft{}Sweat\textquotedblright{}.) Say it once more.
\end{tcolorbox}
\section[échauffer]{échauffer — to warm up; (s'échauffer) to warm up (before sport)}

\label{lesson:FR-C504-echauffer}



\begin{quote}
Before the new one: say the French for to stretch, then the French for sweat.
\end{quote}

\subsection*{You'll want to know: échauffer}
\textbf{échauffer} — \textquotedblleft{}to warm up; (s'échauffer) to warm up (before sport)\textquotedblright{}.

\begin{quote}
Les joueurs s'échauffent avant le match.
\end{quote}

\begin{quote}
\emph{The players warm up before the match.}
\end{quote}

Mostly reflexive in sport: \emph{s'échauffer}. \emph{L'échauffement} — the warm-up.

\begin{cousinweb}
\emph{échauffer} is inherited from Old French \emph{eschaufer} (EtymDB), from Vulgar Latin \emph{excalefāre}, a form of Latin \emph{calefacere} \textquotedblleft{}to make warm\textquotedblright{}, built on \emph{calēre} \textquotedblleft{}to be warm\textquotedblright{} (Lewis \& Short).
\end{cousinweb}

\begin{grammarlens}[title={What you've built}]
One more word for telling, explaining and describing at B1.
\end{grammarlens}

\subsection*{Your turn}
\begin{itemize}
\raggedright
  \item \emph{Say it:} \emph{échauffer}
  \item \emph{Say it:} \emph{échauffer}, once more
  \item \emph{Say it:} a sentence of your own with \emph{échauffer}
  \item \emph{Recall:} say the French for to stretch, then the French for sweat, then say \emph{échauffer} again
\end{itemize}

\begin{tcolorbox}[breakable,colback=teal!4,colframe=teal!35!black,title={Before you move on}]
What does échauffer mean? (\textquotedblleft{}To warm up; (s'échauffer) to warm up (before sport)\textquotedblright{}.) Say it once more.
\end{tcolorbox}
